% 核心结论：地面同时的端点发光在车厢系不同时；光相遇点并非车中点。
% 纯几何时空图，约10x6cm，10pt；L=4,beta=1/3,ctmax=4.2。
% 后端x=ct/3，前端x=4+ct/3；光线A->P<-B，P=(2,2)。
\begin{tikzpicture}[font=\small,>=Latex]
 \def\trainL{4}\def\trainBeta{.3333333333}\def\trainT{4.2}
 \coordinate (A) at (0,0);\coordinate (B) at (\trainL,0);
 \coordinate (P) at ({\trainL/2},{\trainL/2});
 \draw[axis] (-.35,0)--(6.6,0) node[right] {$x$};
 \draw[axis] (0,0)--(0,4.9) node[above] {$ct$};
 \draw[TextGray,thick] (A)--({\trainBeta*\trainT},\trainT) node[above=6pt] {后端};
 \draw[TextGray,thick] (B)--({\trainL+\trainBeta*\trainT},\trainT) node[above=6pt] {前端};
 \draw[DeepBlue,very thick] (A)--(P)--(B);
 \draw[vector] (.65,.65)--(1.3,1.3);
 \draw[vector] (3.35,.65)--(2.7,1.3);
 \draw[Blue,dashed] (A)--(5.8,{5.8*\trainBeta});
 \node[right=5pt] at (5.8,{5.8*\trainBeta}) {$t'=0$};
 \fill (P) circle (1.6pt);\node[above=5pt] at (P) {P};
 \fill (A) circle (1.4pt);\fill (B) circle (1.4pt);
 \node[below left=3pt] at (A) {A};\node[below right=3pt] at (B) {B};
 \draw[TextGray,<->] (0,-.65)--(\trainL,-.65) node[midway,below=3pt] {$L$};
\end{tikzpicture}
